\section{Bài tập 2: Diode in a Series}

Cho mạch điện như hình bên dưới gồm nguồn điện áp $V_1 = 10\text{ V}$, ba diode $D_1, D_2, D_3$ mắc nối tiếp và điện trở tải $R_1 = 1\text{ k}\Omega$. Xác định giá trị điện áp trên các diode $V_{D1}, V_{D2}, V_{D3}$, và dòng điện $I$ trong mạch. Sau đó, tiến hành mô phỏng PSpice để kiểm tra lại kết quả.

Trong bài tập này, giả thiết sử dụng \textbf{mô hình diode thực tế (practical diode model)} với điện áp phân cực thuận ngưỡng là $V_F = 0.7223\text{ V}$.

\given{ex2_given}{0.65\textwidth}{Mạch ba Diode mắc nối tiếp (Hình 2.15)}

\subsection{Tính toán lý thuyết}
Do cực dương nguồn $V_1 = 10\text{ V}$ nối vào Anode của $D_1$ và giá trị nguồn lớn hơn tổng ngưỡng sụt áp thuận của ba diode ($10\text{ V} > 3 \times 0.7223\text{ V} = 2.1669\text{ V}$), cả ba diode $D_1, D_2, D_3$ đều được phân cực thuận và dẫn điện.

\noindent Theo mô hình diode thực tế với $V_F = 0.7223\text{ V}$:
\begin{align*}
  V_{D1} &= 0.7223\text{ V} \\
  V_{D2} &= 0.7223\text{ V} \\
  V_{D3} &= 0.7223\text{ V}
\end{align*}

\noindent Áp dụng định luật KVL quanh vòng kín:
\begin{equation}
  V_1 - V_{D1} - V_{D2} - V_{D3} - V_{R1} = 0
\end{equation}

\noindent Điện áp rơi trên điện trở $R_1$:
\begin{align*}
  V_{R1} &= V_1 - (V_{D1} + V_{D2} + V_{D3}) \\
         &= 10\text{ V} - 3 \times 0.7223\text{ V} = 7.8331\text{ V}
\end{align*}

\noindent Dòng điện $I$ chạy qua mạch nối tiếp:
\begin{equation}
  I = \frac{V_{R1}}{R_1} = \frac{7.8331\text{ V}}{1000\,\Omega} = 7.8331\text{ mA}
\end{equation}

\begin{dapsobox}
  \begin{itemize}[leftmargin=1.5em]
    \item $V_{D1} = V_{D2} = V_{D3} = 0.7223\text{ V}$
    \item $V_{R1} = 7.8331\text{ V}$
    \item $I = 7.8331\text{ mA}$
  \end{itemize}
\end{dapsobox}

\subsection{Mô phỏng PSpice}
\begin{itemize}[leftmargin=1.5em]
  \item Dùng hồ sơ \textbf{Bias Point}, bật hiển thị điện áp và dòng điện; dùng diode \texttt{Dbreak} cho cả ba diode và đối chiếu $V_D$ với $V_F = 0.7223\text{ V}$.
\end{itemize}

\simresult
\simfig{ex2_sim1}{Kết quả mô phỏng Bias Point của Bài tập 2}

\subsection{So sánh và nhận xét}
\begin{table}[H]
  \centering
  \caption{So sánh kết quả tính toán lý thuyết và mô phỏng PSpice cho Bài tập 2}
  \begin{tabular}{lccccc}
    \toprule
    \textbf{Phương pháp} & $V_{D1}\text{ (V)}$ & $V_{D2}\text{ (V)}$ & $V_{D3}\text{ (V)}$ & $V_{R1}\text{ (V)}$ & $I\text{ (mA)}$ \\
    \midrule
    Lý thuyết (Theory)  & $0.7223$ & $0.7223$ & $0.7223$ & $7.8331$ & $7.8331$ \\
    Mô phỏng (PSpice)   & $0.7093$ & $0.7092$ & $0.7093$ & $7.8722$ & $7.8722$ \\
    \bottomrule
  \end{tabular}
\end{table}

\begin{notebox}{Giải thích và ứng dụng thực tế}
\textbf{Giải thích:}
\begin{itemize}[leftmargin=1.5em,itemsep=2pt,topsep=2pt]
  \item Kết quả mô phỏng ($V_D \approx 0.709\text{ V}$, $I = 7.872\text{ mA}$) gần với lý thuyết ($V_F = 0.7223\text{ V}$, $I = 7.833\text{ mA}$), chỉ lệch khoảng $0.013\text{ V}$ trên mỗi diode.
  \item Ba diode giống nhau, mắc nối tiếp nên dòng điện bằng nhau và điện áp phân bố đều, mỗi diode chịu khoảng $0.71\text{ V}$; điện áp còn lại rơi trên $R_1$.
\end{itemize}
\textbf{Ứng dụng thực tế:}
\begin{itemize}[leftmargin=1.5em,itemsep=2pt,topsep=2pt]
  \item Mạch dùng sụt áp thuận gần như cố định của diode để hạ áp nguồn một cách đơn giản, rẻ tiền. Ví dụ module SIM gửi SMS có điện áp cấp tốt khoảng $4.3\text{ V}$: chỉ cần một diode nối tiếp từ nguồn $5\text{ V}$ phổ biến, $5 - 0.7 = 4.3\text{ V}$.
  \item Diode đồng thời bảo vệ module khỏi dòng điện ngược.
\end{itemize}
\end{notebox}
